
\chapter{Conclusion}

This project presented the design and implementation of MetaGPT, a production-ready agentic system for transforming natural language prompts into complete software projects. By modeling the roles of a Manager, Architect, Engineer, and QA agent, and by defining clear SOPs for each, the system brings structure and traceability to LLM-based code generation. The FastAPI backend, LangGraph-based pipeline, and Next.js frontend together provide an integrated environment where users can create projects, observe pipeline execution, inspect generated artifacts, and iteratively refine outputs through chat.

Through implementation and testing, MetaGPT demonstrated that an SOP-driven, multi-agent approach can produce more consistent architectures and clearer reasoning trails than one-shot code generation tools. The system highlights the importance of modular design, explicit state management, and rich visualization when building LLM-powered developer tools. While some limitations remain, particularly around handling highly ambiguous requirements and domain-specific constraints, the work establishes a solid foundation for future research and practical deployment of agentic software engineering systems.
